% New English translation from checked cote 25, site 73 (PDF page 74).
% AI-translated by GPT-6 Astra (OpenAI, ChatGPT Pro mode).
\scanpage{25/73}
\mspage{68}

relative to the base scheme S: V$^1$- V" is indeed\var{\xout{xx}}{}{typewriter strike-out}{ega25_25_073-v0002} smooth over S
(being smooth over X) of relative dimension equal to \uline{one} less than that of
Grass$^m$ over S (as we see by putting k=1 in the formula above).
Thus the results \unsure{of} loc cit apply; in particular we find
that the set of supersingular points of $\hand{\phi}^m$ rel to Grass$^m$
is precisely V" \hand{$\cup$} V$^2$, where V$^2$ is the ramification sub\var{\inserted[typed]{pre}}{pre}{insertion}{ega25_25_073-v0003}scheme
of V$^1$-V" \hand{$\longrightarrow$} G\unsure{r}ass$^m$. Thus we may also say that, outside V",
the supersingular zeros result from the collapse of (at least) two ordinary
singular zeros (but we are not obliged to say so\hand{)}).

\var{\margin{$|$}}{\placedmargin{$|$}}{marginal note placed}{ega25_25_073-v0004} In this way, we essentially have the equivalent of 8.7 a) and b). It \var{\margin{$+$}}{\placedmargin{$+$}}{marginal note placed}{ega25_25_073-v0005}
should be possible to give an equivalent for 8.7 c), \var{\xout{xxx}}{}{typewriter strike-out}{ega25_25_073-v0006} using
the explicit \var{\xout{structure}}{}{typewriter strike-out}{ega25_25_073-v0007} description of the tangent bundle to Grass$^m$ (analogous to
that in the case m=1); \var{\deleted{and which}}{}{deletion}{ega25_25_073-v0008} \var{\inserted{this}}{this}{insertion}{ega25_25_073-v0009} \var{\inserted{also}}{also}{insertion}{ega25_25_073-v0010} implies that for a \var{\inserted[typed]{geometric}}{geometric}{insertion}{ega25_25_073-v0011} point of V$^1$-V" unramified
over Grass$^m$, knowledge of its image in Grass$^m$
implies knowledge of its image in P, provided the first image is a
smooth point of the closed image \var{\inserted[typed]{of}}{of}{insertion}{ega25_25_073-v0012} V$^1$ in Grass$^m$ (we assume S = the spectrum of
a field). I can provide a \var{\xout{slightly}}{}{typewriter strike-out}{ega25_25_073-v0013} more precise statement on request.

Once this point is established, we have obvious corollaries generalizing
8.8, 8.9, 8.11. It is probably just as possible to state
the other propositions of par.8 for arbitrary m. If this requires further
work in writing them up, we can forgo \unsure{this} generalization\hand{,}
\var{\xout{even}}{}{typewriter strike-out}{ega25_25_073-v0014} even if we include the preceding differential developments.
The same holds for the results of n$^o$9.

As for n$^o$10, the situation studied there generalizes to
arbitrary m in the following way. Fix a linear subvariety
C of P of codimension m+1, and consider the projective space Q of linear
\var{\over{\unsure{subvarieties}}{\unsure{subvarieties}}}{\unsure{subvarieties}}{writing over}{ega25_25_073-v0015} L of P of codimension m passing through C. Q is a
\tnote{25-073-1}{The e/g uncertainty in de and the doubtful ending tte in cette are carried by of and this. The doubtful overwritten ariét in sous-variétés is represented by a doubt on the full English word in both branches. The source uses a double quotation mark after V in several places; those marks are retained.}

\var{\margin{73}}{\placedmargin{73}}{marginal note placed}{ega25_25_073-v0016}
